\documentclass[11pt,a4paper]{article}
\usepackage{gtu-el-yazisi}
\begin{document}
\notbilgi{Sayısal İşaret İşleme}{Ders 9: Filtre yapıları (structures)}{Dsp-ders-9.pdf}

\sayfa{1}
\section{Dsp ders 9}
kitapta 6. bölümdeyiz

\section{Structures for discrete-time system}
\begin{enumerate}[label=\arabic*-]
\item Direct form I
\item Direct form II
\item Cascade
\item Parallel structures
\item Transposed form (flow graph reversal)
\end{enumerate}

Önceden filtreleri ikiye ayırmıştık
\elcizimi[0.85]{sekiller/Dsp-ders-9-s1-1.png}[1]
\elcizimi[0.35]{sekiller/Dsp-ders-9-s1-2.png}[1]
\[ y[n] = x[n] * h[n] \]
bilgisayarla yapılabilir \quad $\uparrow$ FIR
\[ \Rightarrow\ y[n] = x[n] * h[n] \ \leftarrow \text{IIR} \]
bilgisayarda yapabilmek için belli bir noktadan sonra kesmek gerekiyor o zaman aslında FIR olur böylece yaklaşıklık yapılır.
Ama biz tam olarak bulmak istiyoruz, o zaman başka bir yöntem kullanacağız.

\section{IIR Filter Design}
\begin{enumerate}[label=\arabic*-]
\item specs. (Teknik özellikler)
\item $h[n]$
\item realize (gerçeklemek) \hfill (this week)
\item implementation (gerçekleştirmek)
\end{enumerate}

\sayfa{2}
\[ y[n] - \sum_{k=1}^{p} a_k\, y[n-k] = \sum_{k=0}^{q} b_k\, x[n-k] \]
$a_k, b_k$ real
\[ H(z) = \frac{\sum_{k=0}^{q} b_k z^{-k}}{1 - \sum_{k=1}^{p} a_k z^{-k}} \qquad \text{rational} \]
Özel sadeleşme durumu \okunamadi{olmazsak} bu IIR bir filtrenin sistem fonksiyonu

\section{Flow notation}
\elcizimi[0.45]{sekiller/Dsp-ders-9-s2-1.png}[2]

\section{Block diagram notation}
\elcizimi[0.55]{sekiller/Dsp-ders-9-s2-2.png}[2]

\[ y[n] = a\,y[n-1] + x[n] \]
\elcizimi[0.9]{sekiller/Dsp-ders-9-s2-3.png}[2]
\[ y[n] = \sum_{k=1}^{p} a_k\, y[n-k] + \underbrace{\sum_{k=0}^{q} b_k\, x[n-k]}_{v[n]} \]

\subsection{Direct Form-I}
\begin{align*}
y[n] &= \sum_{k=1}^{p} a_k\, y[n-k] + v[n]\\
&= a_1 y[n-1] + \cdots + a_p y[n-p] + v[n]
\end{align*}
\[ H(z) = \frac{\sum_{k=0}^{q} b_k z^{-k}}{1 - \sum_{k=1}^{p} a_k z^{-k}} \qquad \text{rational} \]

\sayfa{3}
Önceki girişlerden $q$ tane değer saklanmalı

Önceki çıkışlardan $p$ tane değer saklanmalı
\elcizimi[0.8]{sekiller/Dsp-ders-9-s3-1.png}[3]
\[ H(z) = \frac{H_1(z)}{H_2(z)} = \big(H_1(z)\big)\cdot\left(\frac{1}{H_2(z)}\right) \]
\elcizimi[0.85]{sekiller/Dsp-ders-9-s3-2.png}[3]
\elcizimi[0.85]{sekiller/Dsp-ders-9-s3-3.png}[3]
artı artı diye yaptıklarımız artık eksi olacak $a$'larda sola geçirirken - oluyor
\[ H(z) = \frac{b_0 + b_1 z^{-1} + b_2 z^{-2}}{1 - a_1 z^{-1} - a_2 z^{-2}} \]
\[ y[n] = a_1 y[n-1] + a_2 y[n-2] + b_0 x[n] + b_1 x[n-1] + b_2 x[n-2] \]
\[ y[n] - a_1 y[n-1] - a_2 y[n-2] = b_0 x[n] + b_1 x[n-1] + b_2 x[n-2] \]
önceden bunlar artı idi genel gösterimizde

\sayfa{4}
\[ H(z) = \frac{b_0 + b_1 z^{-1} + b_2 z^{-2}}{1 - a_1 z^{-1} - a_2 z^{-2}} \]
\elcizimi[0.6]{sekiller/Dsp-ders-9-s4-1.png}[4]
\elcizimi[0.65]{sekiller/Dsp-ders-9-s4-2.png}[4]
\elcizimi[0.75]{sekiller/Dsp-ders-9-s4-3.png}[4]

\begin{ornek}
\[ H(z) = \frac{1 + 2z^{-1}}{1 - 1{,}5z^{-1} + 0{,}9z^{-2}} \]
\elcizimi[0.85]{sekiller/Dsp-ders-9-s4-4.png}[4]
\elcizimi[0.65]{sekiller/Dsp-ders-9-s4-5.png}[4]
\end{ornek}

\sayfa{5}
\section{Cascade Structures}
\[ H(z) = A\,\frac{\overbrace{\prod_k (1 - e_k z^{-1})}^{\text{Reel sıfırlar}}\ \prod_k (1 - f_k z^{-1})(1 - f_k^* z^{-1})}{\underbrace{\prod_k (1 - c_k z^{-1})}_{\text{Reel kutuplar}}\ \prod_k (1 - d_k z^{-1})(1 - d_k^* z^{-1})} \]
\begin{align*}
(1 - d_k z^{-1})(1 - d_k^* z^{-1}) &= 1 - 2\,\mathrm{Re}[d_k]\,z^{-1} + |d_k|^2 z^{-2}\\
&= 1 + \beta_{1k} z^{-1} + \beta_{2k} z^{-2}
\end{align*}
O zaman $H(z)$ tekrar bu şekilde yazabiliriz
\[ H(z) = A \prod_k H_k(z) \]
\[ H_k(z) = \frac{1 + \beta_{1k} z^{-1} + \beta_{2k} z^{-2}}{1 - \alpha_{1k} z^{-1} - \alpha_{2k} z^{-2}} \]
s.o.s $\rightarrow$ section

$\swarrow$ second \qquad $\searrow$ order
\elcizimi[0.4]{sekiller/Dsp-ders-9-s5-1.png}[5]
\elcizimi[0.75]{sekiller/Dsp-ders-9-s5-2.png}[5]
SOS Form
\elcizimi[0.75]{sekiller/Dsp-ders-9-s5-3.png}[5]

\begin{ornek}
\begin{align*}
H(z) &= \frac{1 - z^{-1}}{\left(1 + \frac{1}{8}z^{-1}\right)\left(1 + z^{-1} + \frac{1}{4}z^{-2}\right)}\\
&= \left(\frac{1}{1 + \frac{1}{8}z^{-1}}\right)\left(\frac{1 - z^{-1}}{1 + z^{-1} + \frac{1}{4}z^{-2}}\right)\\
&= \left(\frac{1 - z^{-1}}{1 + \frac{1}{8}z^{-1}}\right)\left(\frac{1}{1 + z^{-1} + \frac{1}{4}z^{-2}}\right)
\end{align*}
\end{ornek}

\sayfa{6}
\elcizimi[0.75]{sekiller/Dsp-ders-9-s6-1.png}[6]
Burada 3 tane değeri hafızada saklamak gerekir. (gecikmelerden dolayı)

\section{Parallel implementation}
S.O.S'ların toplamı şeklinde ifade edilmelidir.
\[ H(z) = \frac{\sum_{k=0}^{q} b_k z^{-k}}{1 - \sum_{k=1}^{p} a_k z^{-k}} \]
\[ = \sum_k A_k z^{-k} + \sum_k \frac{B_k}{1 - g_k z^{-1}} + \underbrace{\sum_k \frac{C_k + D_k z^{-1}}{1 - h_{1k} z^{-1} - h_{2k} z^{-2}}}_{\text{en genel hal}} \]
\[ = \sum_k H_k(z) \]
\elcizimi[0.65]{sekiller/Dsp-ders-9-s6-2.png}[6]

\sayfa{7}
\[ H(z) = \underbrace{2z^{-1}}_{H_1(z)} + \underbrace{\frac{1 - z^{-1}}{1 + \frac{1}{2}z^{-1} + \frac{1}{4}z^{-2}}}_{H_2(z)} \]
\elcizimi[0.5]{sekiller/Dsp-ders-9-s7-1.png}[7]
\elcizimi[0.6]{sekiller/Dsp-ders-9-s7-2.png}[7]

\section{Transpozition theorem (Flow-graph reversal theorem)}
Teoreme göre $x[n]$ ve $y[n]$'in yeri değiştirilip tüm okların yönü terslenirse yeni sistem eskinin eşidir.
\[ H_k(z) = \frac{1 + \beta_{1k} z^{-1} + \beta_{2k} z^{-2}}{1 - \alpha_{1k} z^{-1} - \alpha_{2k} z^{-2}} \]
\elcizimi[0.6]{sekiller/Dsp-ders-9-s7-3.png}[7]
\elcizimi[0.6]{sekiller/Dsp-ders-9-s7-4.png}[7]
\begin{align*}
v[n] &= x[n]\,\beta_{1k}\\
&\quad + y[n]\,\alpha_{1k}\\
&\quad + y[n-1]\,\alpha_{2k}\\
&\quad + x[n-1]\,\beta_{2k}\\
y[n] &= x[n] + v[n-1]\\
&= x[n] + x[n-1]\,\beta_{1k} + x[n-2]\,\beta_{2k} + y[n-1]\,\alpha_{1k}\\
&\quad + y[n-2]\,\alpha_{2k}
\end{align*}
\[ \frac{Y(z)}{X(z)} = \frac{1 + \beta_{1k} z^{-1} + \beta_{2k} z^{-2}}{1 - \alpha_{1k} z^{-1} - \alpha_{2k} z^{-2}} \]

\sayfa{8}
Simetriğini alırsak bizim alıştığımız forma geçer
\elcizimi[0.65]{sekiller/Dsp-ders-9-s8-1.png}[8]
\elcizimi[0.75]{sekiller/Dsp-ders-9-s8-2.png}[8]

\section{FIR realization}
$h[n]$: finite
\[ H(z) = \sum_{k=0}^{N} h[k]\, z^{-k} \]
\begin{align*}
\Rightarrow\ y[n] &= \sum_{k=0}^{N} h[k]\, x[n-k] = h[0]x[n] + h[1]x[n-1] + h[2]x[n-2]\\
&\quad + \cdots + h[N]x[n-N]
\end{align*}
\elcizimi[0.9]{sekiller/Dsp-ders-9-s8-3.png}[8]
use transposition theorem:
\elcizimi[0.85]{sekiller/Dsp-ders-9-s8-4.png}[8]

\sayfa{9}
\section{In cascade Structure:}
\[ H(z) = \sum_{k=0}^{N} h[k]\, z^{-k} = \prod_{k=1}^{N} \left(b_{0k} + b_{1k} z^{-1} + b_{2k} z^{-2}\right) \]
\elcizimi[0.8]{sekiller/Dsp-ders-9-s9-1.png}[9]
\[ H(e^{j\omega}) \qquad |H(e^{j\omega})| \qquad \angle H(e^{j\omega}) \]
\[ \kutula{\tau(\omega) = -\frac{d}{d\omega}\,\angle H(e^{j\omega})} \qquad \text{group delay} \]
$\tau(\omega)$ = constant
\elcizimi[0.45]{sekiller/Dsp-ders-9-s9-2.png}[9]
\elcizimi[0.5]{sekiller/Dsp-ders-9-s9-3.png}[9]
Önceden böyle bir örnek göstermiştik

\fbox{$0{,}2\pi$}

$0{,}4\pi$

$0{,}8\pi$
\[ \kutula{\angle H(e^{j\omega}) = \alpha\omega + \beta} \quad \text{yapısında ise} \quad \tau(\omega) = \text{constant} \]
Linear Phase Filter

IIR $\times$ (hiç bir zaman sağlayamaz) \qquad FIR sağlayabilir

\sayfa{10}
\subsection{Type-I}
\[ h[n] = h[M-n] \qquad 0 \le n \le M \]
$M$: even integer (çift)
\[ M = 4 \;\Rightarrow\; 0 \le n \le 4 \]
\begin{align*}
h[0] &= h[4]\\
h[1] &= h[3]\\
h[2]&
\end{align*}
\elcizimi[0.45]{sekiller/Dsp-ders-9-s10-1.png}[10]

\subsection{Type II}
\[ h[n] = h[M-n] \qquad 0 \le n \le M \]
$M$: odd integer (tek sayı)
\[ M = 5 \;\Rightarrow\; 0 \le n \le 5 \]
\begin{align*}
h[0] &= h[5]\\
h[1] &= h[4]\\
h[2] &= h[3]
\end{align*}
\elcizimi[0.65]{sekiller/Dsp-ders-9-s10-2.png}[10]

\subsection{Type III}
\[ h[n] = -h[M-n] \qquad 0 \le n \le M \]
$M$: even integer (çift sayı)

$M = 4$
\elcizimi[0.5]{sekiller/Dsp-ders-9-s10-3.png}[10]

\subsection{Type IV}
\[ h[n] = -h[M-n] \qquad 0 \le n \le M \qquad M\text{: odd integer} \]
$M = 5$
\elcizimi[0.6]{sekiller/Dsp-ders-9-s10-4.png}[10]

\end{document}
